\section{Uncertainty, robustness and limitations}\label{sec:discussion}

\subsection{Four layers of uncertainty}

The intervals reported in Sections~\ref{sec:q3} and~\ref{sec:q4} are the narrowest of the four
layers of uncertainty that the analysis contains. Reporting them alone would give a false sense
of precision.

\begin{table}[htbp]
  \centering
  \caption{Four layers of uncertainty, their typical magnitude, and how they were handled.}\label{tab:uncertainty}
  \small
  \begin{tabular}{p{2.6cm}p{4.6cm}p{5.6cm}}
    \toprule
    Layer & Magnitude & Treatment \\
    \midrule
    Sampling (blocks) & two-voice radius $\pm 5$ m by block bootstrap & block bootstrap over 30 blocks, not rows \\
    Model family & two-voice radius 158--203 m (logistic vs empirical) & report the range, do not choose \\
    Threshold and weights & radius 203 $\to$ 143 m when $\tau$ moves 0.5 $\to$ 0.7; 12--67 m under envelope weight 0.5 & full sensitivity table, Section~\ref{sec:q4-sens} \\
    Definition of outcomes & any response (all 278 rows) 108 $\to$ 112 across S1--S3 & definitional sensitivity, Table~\ref{tab:defsens} \\
    \bottomrule
  \end{tabular}
\end{table}

The ordering is instructive: the sampling layer, the only one that conventional statistics
reports, is smaller than the others by an order of magnitude. A paper that reported
$203.4\pm5$ m and stopped would be technically correct and practically misleading.

\subsection{What would change the conclusions}

Three statements survive every sensitivity analysis and three do not.

\emph{Robust.} (i) Natural calls elicit far more responses than controls, by a factor between 13
and 37 depending on the matching used. (ii) Response probability declines with distance and with
group size. (iii) The distribution of response combinations differs between call and control.

\emph{Fragile.} (i) The ordering of call types at the 0.5 threshold is model-dependent: a
four-parameter Weibull places modulated marginally ahead of two-voice, while the empirical curve,
the GAM and the three-parameter Weibull place two-voice ahead. (ii) The claim that two-voice
calls are most robust at long range holds for the spectrogram and amplitude-spectrum metrics and
reverses for the envelope metric at 256 m. (iii) The apparent effect of context disappears after
adjustment for group size, and the data cannot say whether the remaining association with group
size is behavioural or observational.

\subsection{Limitations}

\textbf{The observation window is not recorded (L1).} We assume $W=60$ s because it is the upper
bound of the duration field and exceeds the largest observed latency; the sensitivity over 48,
90 and 120 s bounds the effect, and if the true window were shorter some non-responders would be
misclassified. \textbf{Nine zero latencies are excluded from the latency model (L2)}; they are
immediate responses at the resolution of the protocol and would form a point mass that a
continuous likelihood cannot accommodate. \textbf{Playback and propagation trials cannot be
linked (L3)}, because no recording identifier connects them, so $q(d)$ in
Equation~\eqref{eq:decomp} must be imported rather than estimated. \textbf{Date is confounded
with call type (L4)}: the association is strong (Cram\'er's $V=0.528$), and a date-fixed-effects
fit raises the modulated and two-voice odds ratios to about 46 and 39; although the intervals
widen (13--163 and 13--116), both remain significant, so the association survives the date
confound. The estimates are imprecise because date and call type are strongly associated. We keep
the unadjusted model as primary and report the adjusted direction as a caveat. \textbf{The propagation data have a small effective sample (L5)}: thirty
blocks support about 46 effective observations per distance, which the block bootstrap reflects
but cannot enlarge. Finally, \textbf{all conclusions are conditional on a single site and a
single season (L6)} and should not be extrapolated without re-estimation.
\textbf{Recruitment is extrapolated beyond the playback range (L7)}: playback distances stop at
200 m, whereas the CEI curve is evaluated to 256 m, so the two-voice radius of 203 m sits at the
edge of the observed range and should be read as an upper bound of about 200 m. \textbf{The CEI
bands reflect fidelity sampling only (L8)}: the block bootstrap resamples the 30 propagation
blocks but does not propagate the uncertainty in the GEE coefficients for $P_{\mathrm{recruit}}$,
so the bands in Figure~\ref{fig:q4-cei} understate the total uncertainty.

\subsection{Strengths and weaknesses of the model}

The strength of the approach is that it uses the design the data actually have: the pairing is
exploited rather than discarded, censoring is modelled rather than deleted, the block structure
of the propagation data is respected in the CEI intervals of Section~\ref{sec:q4}, and the index
is built from two separately estimated components whose provenance is traceable.

The weakness is that the two components are measured on different scales. Fidelity is a
correlation between waveforms, recruitment is a probability of behaviour, and their product is
dimensionless but not calibrated to any ecological outcome. The index orders distances and call
types correctly; it does not tell us how many calories a bird saves, or how many recruits a
colony needs. Closing that gap would require a fitness-linked experiment that the present data
cannot supply.
